\centerline{\bf \Huge Комби профилактика}

\begin{flushright}
        \scriptsize{}
\end{flushright}


\problem{1}{На доске написано несколько натуральных чисел. Сумма любых двух из них является натуральной степенью двойки. Какое наибольшее количество различных чисел может быть написано на доске?}

% Источник: Турнир городов, 2012/13, весенний тур, сложный вариант,
% 8--9 классы, задача 1.
% https://problems.ru/view_problem_details_new.php?id=116683

\medskip

\problem{2}{Петя заметил, что у всех его $25$ одноклассников различное количество друзей в этом классе. Сколько друзей может быть у самого Пети? Укажите все возможные ответы.}

% Источник: Турнир городов, 1992/93, весенний тур, основной вариант,
% 8--9 классы; автор С.\,И.~Токарев.
% https://problems.ru/view_problem_details_new.php?id=98170

\medskip

\problem{3}{На окружности расставлено несколько положительных чисел, каждое из которых не превосходит~$1$. Докажите, что окружность можно разделить на три дуги так, чтобы суммы чисел на любых двух соседних дугах отличались не более чем на~$1$. Если на дуге нет чисел, сумма на ней считается равной нулю.}

% Источник: Турнир городов, 2005/06, осенний тур, основной вариант,
% 10--11 классы, задача 4; автор М.\,И.~Малкин.
% https://problems.ru/view_problem_details_new.php?id=65833

\medskip

\problem{4}{В шахматном турнире участвовали $11$ человек: четверо россиян и семеро иностранцев. Каждые два участника сыграли между собой две партии. За победу начислялось $1$ очко, за ничью~--- $\frac12$ очка, за поражение~--- $0$ очков. Все участники набрали различное количество очков, причём сумма очков российских шахматистов оказалась равна сумме очков иностранных шахматистов. Могло ли случиться, что среди трёх участников, занявших первые места, не было ни одного россиянина?}

% Источник: Московская математическая регата, 2014/15,
% 10 класс, задача 10.3.3.
% https://problems.ru/view_problem_details_new.php?id=65178

\medskip

\problem{5}{На окружности отмечена $21$ точка. Каждая пара отмеченных точек является концами двух дуг окружности. Докажите, что среди всех этих дуг найдётся не менее ста дуг, угловая мера которых не превосходит $120^\circ$.}

% Источник: Турнир городов, 1986/87, осенний тур,
% 9--10 классы, задача 7; автор А.\,Ф.~Сидоренко.
% Также задача М1055 журнала ``Квант''.
% https://problems.ru/view_problem_details_new.php?id=97920

\medskip

\problem{6}{В стране есть две столицы: Южная и Северная, а также несколько других городов. Некоторые города соединены односторонними дорогами. Часть дорог является платной. Известно, что на любом пути из Южной столицы в Северную встречается не менее десяти платных дорог. Докажите, что платные дороги можно распределить между десятью компаниями так, чтобы на любом пути из Южной столицы в Северную встречалась дорога каждой из десяти компаний.}

% Источник: Турнир городов, 2009/10, осенний тур, основной вариант,
% 10--11 классы, задача 5; авторы И.\,В.~Нетай, Д.\,В.~Баранов.
% Формулировка уточнена указанием, что дороги односторонние.
% https://problems.ru/view_problem_details_new.php?id=116247

\medskip

\problem{7}{В каждой клетке квадратной таблицы написано число. Известно, что в каждой строке сумма двух наибольших чисел равна~$a$, а в каждом столбце сумма двух наибольших чисел равна~$b$. Докажите, что $a=b$.}

% Источник: Турнир городов, 2010/11, весенний тур, сложный вариант,
% 10--11 классы, задача 6.
% Также: Московская математическая олимпиада, 2011, 8 класс, задача 6.
% https://problems.ru/view_problem_details_new.php?id=116213

\medskip

\problem{8}{В Академии наук работают $999$ академиков. Каждая научная тема интересует ровно троих академиков, а у каждых двух академиков есть ровно одна тема, интересная им обоим. Докажите, что можно выбрать $250$ тем так, чтобы каждый академик интересовался не более чем одной из выбранных тем.}

% Источник: заключительный этап Всероссийской олимпиады школьников,
% 2010/11, 11 класс, задача 3; автор А.~Магазинов.
% https://problems.ru/view_problem_details_new.php?id=116648